\section{Conclusion}\label{sec:conclusion}

This paper identifies a structural lever in NPU intra-kernel scheduling: the
$2\times$ asymmetry between clean and dirty spill costs means that schedule order
controls not only peak live bytes, but also the cost composition of the buffers
available for eviction in capacity-pressure windows. When spills are unavoidable,
keeping more cheaply evictable clean reserve substantially reduces real off-chip
traffic. Around this finding we give a spill-inevitability certificate
(Theorem~\ref{thm:spill-inev}), a conditional overflow-area approximation
(Theorem~\ref{thm:area-extra}), and a Belady-margin stability result
(Proposition~\ref{prop:belady-margin}) that explains the empirical pattern of an
insensitive victim rule alongside a highly sensitive schedule order.

Experiments support the claim at four levels. On the public benchmark cases,
clean/dirty-blind memory-pressure schedulers pay $2.4\text{--}26\times$ more P2
spill traffic than our method (Section~\ref{sec:exp-baselines}). The synthetic
distribution shows that the gains concentrate in the capacity-bound region
delimited by the certificate; in order-reachable instances a simple order already
avoids spills and our method has no systematic advantage
(Section~\ref{sec:exp-mechanism}). Small-graph CP-SAT oracles confirm fixed-order
engine optimality and the empirical constant ranges of the two bounds, and the
controlled ablation isolates the $2\times$ clean/dirty cost gap, showing that the
benefit is not an artifact of particular instances.

Future work includes extending spill-cost-aware liveness shaping to multi-core
scheduling and cross-core cache coherence; supporting graphs with dynamic control
flow such as conditional branches and loops; and modeling cross-level spill
cascades in deeper memory hierarchies.
